\documentclass[10pt,a4paper]{article}
\usepackage{fontspec,amsmath,amssymb,graphicx}
\usepackage[margin=18mm]{geometry}
\setmainfont{Times New Roman}
\setlength{\parindent}{0pt}
\setlength{\parskip}{5pt}
\allowdisplaybreaks
\sloppy
\emergencystretch=2em
\begin{document}
\section*{VD\_4\_1\_2\_02 — Lời giải bổ sung}
\subsection*{Kiến thức và trình tự giải}
\textbf{Bước 1 — Đọc hệ số và điều kiện biên.} Viết phương trình dưới dạng $u_t=Ku_{xx}$, $K>0$. Nếu đề ký hiệu tốc độ $a$ và viết $u_t=a^2u_{xx}$ thì thay $K=a^2$; không nhầm $a$ với $K$. Kiểm tra điều kiện đầu $f$ có khớp giá trị biên tại $t=0$ hay không. Sự không tương thích tại góc không ngăn nghiệm trơn với $t>0$, nhưng ngăn tính liên tục đến góc khi biên Dirichlet khác vết của $f$.

\textbf{Bước 2 — Tách biến và giải bài toán không gian.} Thế $u=XT$ vào PDE, chia cho $KXT$ để có $X''/X=T'/(KT)=-\lambda$. Do đó $X''+\lambda X=0$, $T'+K\lambda T=0$. Với $\lambda=k^2>0$, $X=A\cos kx+B\sin kx$; thay từng điều kiện biên để tìm $k$. Hai biên Dirichlet cho $X_n=\sin(n\pi x/L)$, $n\ge1$; hai biên Neumann cho $1$ và $\cos(n\pi x/L)$; biên Dirichlet–Neumann cho sine, biên Neumann–Dirichlet cho cosine, với $k_n=(n+1/2)\pi/L$, $n\ge0$. Trường hợp $\lambda=0$ phải giải riêng $X=A+Bx$; chỉ giữ mode hằng nếu biên cho phép.

\textbf{Bước 3 — Tính hệ số, không chỉ ghi tích phân.} Đặt $N_n=\int_0^L X_n^2dx$, $c_n=N_n^{-1}\int_0^L fX_n\,dx$. Với các mode sine/cosine chuẩn, $N_n=L/2$; mode hằng có $N_0=L$. Nếu $f$ là đa thức, tích phân từng phần đến hết; nếu $f$ chia đoạn, tách đúng cận; nếu $f$ lượng giác, dùng tích–tổng. Chẳng hạn với $k=n\pi/L$, $f(0)=f(L)=0$,


\[
\int_0^Lf\sin kx\,dx=\frac1k\int_0^Lf'\cos kx\,dx
=-\frac1{k^2}\int_0^Lf''\sin kx\,dx.
\]


Các hạng biên phải được tính trước khi bỏ. Khi tích–tổng tạo mẫu $k-m=0$, tính trực tiếp tích phân ở tần số đó.

\textbf{Bước 4 — Giải ODE và ghép nghiệm.} Nhân $T_n'+K\lambda_nT_n=0$ với $e^{K\lambda_nt}$, được $(e^{K\lambda_nt}T_n)'=0$. Vì $T_n(0)=c_n$,


\[
u(x,t)=\sum_n c_ne^{-K\lambda_nt}X_n(x).
\]


Thêm $c_0=L^{-1}\int_0^Lf$ cho Neumann hai đầu. Kiểm tra từng mode: $X_n''=-\lambda_nX_n$, nên đạo hàm theo $t$ và $K$ lần đạo hàm theo $x$ cấp hai bằng nhau.

\textbf{Bước 5 — Kiểm tra dữ liệu, hội tụ và thời gian dài.} Với $t\ge\varepsilon>0$, thừa số $e^{-K\lambda_nt}$ cho phép lấy đạo hàm chuỗi. Với dữ liệu có bước nhảy, giới hạn ban đầu tại bước nhảy là trung bình hai vết. Nếu $f\in L^1$, $|c_n|\le2\|f\|_1/L$ cho hệ chuẩn, nên tổng các mode dương tiến về không đều theo $x$. Neumann giữ trung bình; Dirichlet và biên hỗn hợp có giới hạn không. Một nghiệm giảm theo hàm mũ và khác không không thể tuần hoàn theo thời gian.
\subsection*{Lời giải áp dụng cho bài}
\textbf{Bước 1. Cơ sở Neumann.} $X'(0)=X'(L)=0$ cho mode hằng $1$ và $X_n=\cos(n\pi x/L)$, $n\ge1$. Chuẩn lần lượt $L,L/2$, nên
\[u=c_0+\sum_{n\ge1}c_ne^{-K(n\pi/L)^2t}\cos(n\pi x/L),\quad c_0=\frac1L\int f,\quad c_n=\frac2L\int f\cos(n\pi x/L).\]
\textbf{Bước 2. Dữ liệu (i).} $\cos^2(\pi x)=(1+\cos2\pi x)/2$, nên với $K=L=1$, $c_0=c_2=1/2$, các hệ số khác bằng không:
\[u=\frac12+\frac12e^{-4\pi^2t}\cos2\pi x.\]
Đạo hàm biên là $-\pi e^{-4\pi^2t}\sin2\pi x$, bằng không ở $0,1$. Thế $t=0$ trả lại $\cos^2\pi x$; đạo hàm $u_t,u_{xx}$ đều bằng $-2\pi^2e^{-4\pi^2t}\cos2\pi x$. Nghiệm trơn và tiến đều về $1/2$.

\textbf{Bước 3. Dữ liệu (ii).} $K=2,L=\pi,f=\sin x$, $c_0=2/\pi$. Với $n\ne1$,
\begin{align*}
\int_0^\pi\sin x\cos nx\,dx
&=\frac12\int_0^\pi[\sin((1+n)x)+\sin((1-n)x)]dx\\
&=\frac{1+(-1)^n}{1-n^2}.
\end{align*}
Khi $n=1$, $\int\sin x\cos x=[\sin^2x/2]_0^\pi=0$. Vì thế $c_n=0$ khi $n$ lẻ, $c_{2j}=4/[\pi(1-4j^2)]$. Suy
\[u=\frac2\pi+\frac4\pi\sum_{j\ge1}\frac{e^{-8j^2t}}{1-4j^2}\cos2jx.\]
\textbf{Bước 4. Kiểm tra tính thông thường.} Hệ số $O(j^{-2})$ cho hội tụ đều đến $t=0$, trả lại $f$ liên tục. Khi $t>0$, chuỗi đạo hàm hội tụ đều, đúng PDE và Neumann. Vì $f'(0)=1,f'(\pi)=-1$, đạo hàm không thể liên tục đến hai góc và đồng thời bằng $0$ trên biên. Nghiệm vẫn thuộc lớp $C^{2,1}$ với $t>0$ và liên tục trên bao đóng, là lớp nghiệm thông thường được dùng trong ví dụ; không khẳng định trơn đến góc. Ước lượng
\[\left|u-\frac2\pi\right|\le\frac4\pi e^{-8t}\sum_{j\ge1}\frac1{4j^2-1}\longrightarrow0\]
chứng minh giới hạn đều $2/\pi$.
\begin{center}\includegraphics[width=.97\linewidth,height=11cm,keepaspectratio]{Do_thi/VD_4_1_2_02.png}\end{center}
\textit{Đồ thị minh họa; với nghiệm số, xem phương pháp và tham số trong mã Python kèm theo.}\end{document}